\chapter{Built-in Commands}
\label{ch:commands}

\index{commands}

Acme's built-in commands are capitalized words that you execute by
middle-clicking (\button{2}) on them.  They typically appear in window
tag bars but can be typed and executed anywhere.

\section{Clipboard: Cut, Snarf, Paste}
\label{sec:clipboard}

\index{Cut}
\index{Snarf}
\index{Paste}
\index{snarf buffer}
\index{clipboard}

Acme maintains an internal buffer called the \emph{snarf buffer}.  This
is also synchronized with the system clipboard via SDL3, so you can
copy and paste between acme and other applications.

\begin{description}[style=nextline]
\item[\cmd{Cut}]
Delete the selected text and place it in the snarf buffer.  The
window is marked dirty.

\item[\cmd{Snarf}]
Copy the selected text to the snarf buffer \emph{without} deleting
it.  The window is not marked dirty.

\item[\cmd{Paste}]
Replace the current selection with the contents of the snarf buffer.
If there is no selection, insert at the cursor position.
\end{description}

The mouse chords \chord{1}{2} (Cut), \chord{1}{3} (Paste), and
\chord{1}{2}+\button{3} (Snarf) perform the same operations more
quickly.  See Section~\ref{sec:chords}.

To paste text from another application into acme: copy in the other
application (e.g., \key{Ctrl+C}), then execute \cmd{Paste} in acme or
use the \chord{1}{3} chord.

\section{Undo and Redo}
\label{sec:undo}

\index{Undo}
\index{Redo}

\begin{description}[style=nextline]
\item[\cmd{Undo}]
Undo the last change or group of changes.  Acme groups related
changes (e.g., all characters typed before a pause) into a single
undo step.

\item[\cmd{Redo}]
Reverse the last \cmd{Undo}.
\end{description}

Undo operates on a per-file basis.  If the same file is open in
multiple windows, undoing in one window affects all windows showing that
file.

\section{File Operations: Get, Put, Putall}
\label{sec:fileops}

\index{Get}
\index{Put}
\index{Putall}

\begin{description}[style=nextline]
\item[\cmd{Get}]
Reload the file from disk, replacing the window's contents.  If the
window has unsaved changes, acme warns you (execute \cmd{Get} a
second time to override).

With an argument (\cmd{Get \textit{filename}}), load the named file
into the current window without changing the window's name.

\item[\cmd{Put}]
Write the window's contents to disk.  Uses the filename shown in the
tag bar.

With an argument (\cmd{Put \textit{filename}}), write to the named
file instead (``Save As'').

\item[\cmd{Putall}]
Write all dirty windows to their respective files.  Only saves
windows whose names correspond to real files.  Skips scratch
windows, directory windows, and windows with no filename.
\end{description}

\section{Window Management}
\label{sec:winmgmt}

\index{New}
\index{Del}
\index{Delete}
\index{Zerox}
\index{Newcol}
\index{Delcol}
\index{Sort}

\begin{description}[style=nextline]
\item[\cmd{New}]
Create a new empty window.  With arguments (\cmd{New
\textit{file1 file2...}}), open those files in new windows.

\item[\cmd{Del}]
Close the current window.  If the window has unsaved changes, acme
warns you.  Execute \cmd{Del} a second time to close anyway.

\item[\cmd{Delete}]
Close the current window without checking for unsaved changes.

\item[\cmd{Zerox}]
Clone the current window, creating a second view of the same file.
Both windows show the same underlying file; changes in one appear in
the other.

\item[\cmd{Newcol}]
Create a new column.

\item[\cmd{Delcol}]
Delete the current column and all its windows.  Warns if any window
has unsaved changes.

\item[\cmd{Sort}]
Sort the windows in the current column alphabetically by filename.
\end{description}

\section{Session Management: Dump, Load, Exit}
\label{sec:session}

\index{Dump}
\index{Load}
\index{Exit}

\begin{description}[style=nextline]
\item[\cmd{Dump}]
Save the complete acme session state---window layout, filenames, font
settings, and positions---to a file.  Default: \file{\$HOME/acme.dump}.

With an argument (\cmd{Dump \textit{filename}}), save to the named
file.

\item[\cmd{Load}]
Restore a previously saved session from a dump file.  Default:
\file{\$HOME/acme.dump}.

\item[\cmd{Exit}]
Quit acme.  If any window has unsaved changes, acme warns you.
Execute \cmd{Exit} a second time to quit anyway.
\end{description}

The \cmd{Dump}/\cmd{Load} pair lets you preserve your workspace
across sessions.  It is common to dump before quitting and load when
restarting.

\section{Search: Look}
\label{sec:look}

\index{Look}

\begin{description}[style=nextline]
\item[\cmd{Look}]
Search for text in the window body.  Uses the current selection as
the search term.  If there is no selection, uses the word nearest
the cursor.

With an argument (via the 2-1 chord or by typing \cmd{Look
\textit{text}}), search for that text.

The search is \textbf{literal} (not a regex) and wraps around the
end of the file.  For regex search, use the \cmd{Edit} command
(Chapter~\ref{ch:edit}) or right-click on a \texttt{:/regex/}
address.
\end{description}

\section{Display: Font, Tab, Indent}
\label{sec:display}

\index{Font}
\index{Tab}
\index{Indent}
\index{auto-indent}

\begin{description}[style=nextline]
\item[\cmd{Font}]
Toggle between proportional and fixed-width fonts.

With an argument (\cmd{Font \textit{fontname}}), set the window's
font.  Font names are fontconfig names with a size suffix (e.g.,
\file{DejaVuSans/12a/font}).

\item[\cmd{Tab}]
With no argument: print the current tab width.

With a numeric argument (\cmd{Tab 8}): set the tab display width for
the current window.

\item[\cmd{Indent}]
Control auto-indentation.

\begin{itemize}
\item \cmd{Indent on} / \cmd{Indent off} --- affect the current window.
\item \cmd{Indent ON} / \cmd{Indent OFF} --- affect all windows globally.
\end{itemize}
\end{description}

\section{Process Control: Kill, ID}
\label{sec:process}

\index{Kill}
\index{ID}

\begin{description}[style=nextline]
\item[\cmd{Kill}]
Send a kill signal to running commands.  With arguments, kills
commands whose names match.

\item[\cmd{ID}]
Print the current window's numeric ID.
\end{description}

\section{Include Paths: Incl}
\label{sec:incl}

\index{Incl}

\begin{description}[style=nextline]
\item[\cmd{Incl}]
With no arguments: list the current supplementary include directories.

With arguments (\cmd{Incl \textit{/path/to/headers}}): add
directories to the search path used when right-clicking on
\texttt{<filename.h>} patterns.
\end{description}

The default search path includes \file{/usr/include} and
\file{/usr/local/include}.  The \cmd{Incl} command adds project-specific
directories.

\section{Send}
\label{sec:send}

\index{Send}

\begin{description}[style=nextline]
\item[\cmd{Send}]
In a Win window (see Chapter~\ref{ch:shell}), send the text between
the prompt boundary and the end of the body to the shell.

In a regular window, append the snarf buffer contents to the end
of the body.
\end{description}

\section{Win}
\label{sec:win-brief}

\index{Win}

The \cmd{Win} command opens an interactive shell session in an acme
window.  See Chapter~\ref{ch:shell} for full details.
